\chapter{DirectX 1--3 and free-direct}
\label{ch:directx-time-machine}
\label{ch:dx3-freedirect-renderer}
\label{ch:dx3-freedirect-backend}

Four identities meet here, but only three are Microsoft-era rungs.  \texttt{DIRECTX1},
\texttt{DIRECTX2}, and \texttt{DIRECTX3} use historical DirectDraw/Direct3D interfaces on
Windows.  \texttt{FREEDIRECT} is a portable 2D renderer over the separate \texttt{free-direct}
project and is not an implementation of the \texttt{DIRECTX3} selector.

\begin{longtable}{p{0.16\textwidth}p{0.24\textwidth}p{0.24\textwidth}p{0.22\textwidth}}
\toprule
Identity & Native surface & 3D route & RT / MRT / query \\
\midrule
\endfirsthead
\toprule
Identity & Native surface & 3D route & RT / MRT / query \\
\midrule
\endhead
\texttt{DIRECTX1} & DirectDraw v1 & unavailable; base fallback rejects & off-screen surface / no / no \\
\texttt{DIRECTX2} & DirectDraw and Direct3D v2 & fixed-function \texttt{DrawPrimitive} & yes / no / no \\
\texttt{DIRECTX3} & DirectDraw v2 and Direct3D v2 & fixed-function \texttt{DrawPrimitive} & yes / no / no \\
\texttt{FREEDIRECT} & free-direct surface API & unavailable; base fallback rejects & yes / no / no \\
\bottomrule
\end{longtable}

\section{DIRECTX1: before Direct3D existed}

DirectX~1 shipped DirectDraw but no Direct3D. CNA's renderer is consequently a 2D
implementation, not an unfinished 3D renderer.  Its off-screen surfaces provide
\cnaclass{RenderTarget2D}; MRT and occlusion queries do not exist at this API level.

The extended primitive methods are inherited.  Their base route reaches colored rendering,
and this 2D family rejects that request.  This is the correct visible failure for a permanent
API ceiling.  Reporting successful effect-aware 3D would be more misleading than throwing.

Build and execution are Windows-shaped.  The repository can cross-compile with mingw-w64 and
exercise the binary under Wine's DirectDraw implementation, but that proves a compatibility
layer as well as CNA.  Native historical drivers remain a different environment claim.

\section{DIRECTX2: fixed function without execute buffers}

Direct3D~2 is often associated with execute buffers, yet CNA submits immediate
\texttt{DrawPrimitive} calls. A spike established that Wine's implementation accepts the route
and provides the transform, lighting, texture,
depth, and raster behavior needed by the bounded renderer.

Both extended draw entry points are native fixed-function adapters.  CNA's modern draw packet
is reduced to texture, matrix, lighting, fog, alpha, and state concepts available in the old
pipeline.  Custom programmable effects are outside the API's design.  2D targets work;
multiple attachments and queries do not.

\section{DIRECTX3: a small port, not a renamed free-direct}

This renderer uses the DirectDraw~2 and Direct3D~2 interface generation available with the
DirectX~3 SDK surface.  It landed under a temporary \texttt{DX30} name because the old
\texttt{DIRECTX3} selector was still occupied by what is now \texttt{FREEDIRECT}.  The
identity cleanup separated historical Microsoft DirectX from the portable sibling project.

The measured source delta from the DirectX~2 port was small, roughly 3.7 percent, but that
number describes lineage rather than semantic equivalence.  Interface acquisition,
presentation, and platform behavior still belong to a separate family and build target.
Its draw and target feature shape matches the preceding rung: native fixed-function draws,
one 2D target, no MRT, and no occlusion queries.

\section{Why the first three rungs matter}

The sequence isolates what the API itself added.  DirectX~1 can compose surfaces but cannot
express 3D.  DirectX~2 introduces an immediate fixed-function route sufficient for CNA's
bounded 3D mapping.  DirectX~3 changes the surrounding SDK/interface generation without
turning that fixed pipeline into a programmable one.  The same CNA call surface therefore
reveals a historical API boundary instead of flattening every old renderer into one
generic legacy label.

\section{FREEDIRECT: a separate portable 2D branch}

\texttt{FREEDIRECT} is CNA's second 2D-only renderer, and its odd one out architecturally: it does
not use SDL3's 2D API at all, instead fronting \texttt{IDirectDraw} / \texttt{IDirectDrawSurface}-shaped
calls against \texttt{free-direct}, the narrow DirectDraw reimplementation introduced in
Chapter~\ref{ch:ecosystem-map} and covered as its own project later in this book. This chapter
covers only what is specific to CNA's own renderer layer on top of it; this book's
\texttt{free-direct} chapter covers the library itself.

\section{An unanticipated architectural constraint}
\label{sec:freedirect-shadow-backbuffer}

\texttt{free-direct}'s \texttt{IDirectDrawSurface::Lock()} never exposes a writable
pointer for the \emph{primary} surface---only for offscreen ones. Instead of treating the
primary surface as a render target directly, this renderer owns an internal, always-lockable
``shadow backbuffer'': every \texttt{Clear()} and every \cnaclass{SpriteBatch} draw targets
that shadow surface, and \texttt{Present()} is a single identity \texttt{Blt()} copying it
onto the primary surface. This design makes the renderer possible.

\section{A capability beyond SDL\_RENDERER}

\texttt{TextureAddressMode::Wrap} and \texttt{Mirror} --- the two features
Chapter~\ref{ch:sdl-renderer} named as formally BLOCKED on \texttt{SDL\_RENDERER} --- have
implemented sampling paths here. This renderer's CPU compositor
already samples per-source-pixel for any non-identity draw, so wrap/mirror addressing costs
nothing extra to add. \texttt{FREEDIRECT} also implements separate Opaque, AlphaBlend,
NonPremultiplied, and Additive formulas, in contrast to
the \texttt{SOFTWARE} renderer's single collapsed baseline formula.

\subsection{Address modes beyond SDL\_RENDERER}

The fixture from Chapter~\ref{ch:sdl-renderer} draws a two-texel red/green texture through a
four-texel source rectangle onto four destination pixels. With point filtering, each destination
sample is unambiguous. FREEDIRECT reads back all four red-channel values for each address mode:

\begin{lstlisting}
Texture2D tex(dev, 2, 1);
std::vector<Color> px = {Color(255, 0, 0, 255), Color(0, 255, 0, 255)}; // texel 0 = red, 1 = green

SamplerState point;
point.setFilterProperty(TextureFilter::Point);
point.setAddressUProperty(TextureAddressMode::Wrap); // or Mirror, or Clamp
point.setAddressVProperty(TextureAddressMode::Wrap);

spriteBatch_->Begin(SpriteSortMode::Deferred, BlendState::AlphaBlend, &point, nullptr, nullptr);
spriteBatch_->Draw(tex, Rectangle(0, 0, 4, 1), Rectangle(0, 0, 4, 1), Color::White);
spriteBatch_->End();
// Read back R at destination x = 0,1,2,3:
\end{lstlisting}

\texttt{Wrap} yields \texttt{(255, 0, 255, 0)}, or red/green twice. \texttt{Mirror} yields
\texttt{(255, 0, 0, 255)}, and \texttt{Clamp} yields \texttt{(255, 0, 0, 0)}. SDL Renderer
provides only the clamp-shaped result through this draw path. FREEDIRECT can implement the three
index rules inside its existing per-source-pixel CPU loop; SDL's texture-blit route exposes no
equivalent sampler control.

\section{Corrected FREEDIRECT defects}

The renderer audit corrected several independent faults:

\begin{itemize}
  \item \texttt{Clear()} wrote alpha 255 through its color-fill route; locked-pixel writes now
        preserve all four requested channels.
  \item the general CPU blend path omitted source-alpha multiplication;
  \item \texttt{CreateOcclusionQuery()} threw instead of using the project's null-safe query;
  \item blend-preset detection compared factors but ignored \cnaclass{BlendFunction}, so a
        custom subtract state could be mistaken for a preset; and
  \item display-free tests were unnecessarily registered as requiring X11.
\end{itemize}

The current regressions cover the behavior, while historical test-count bookkeeping is not used
as a capability claim.

\section{RenderTarget2D: the same shadow-backbuffer trick}

\cnaclass{RenderTarget2D} reuses the shadow-backbuffer design once per target. The private
render-target renderer owns an always-lockable
offscreen surface, defined (like \texttt{FreeDirectTextureRenderer}) entirely
inside \texttt{FreeDirectRenderer.cpp} so \texttt{<ddraw.h>} never leaks into a header. Binding
one is a pure redirect: \texttt{Clear()} and \texttt{ReadBackbuffer()} both route through an
internal \texttt{Impl::ActiveSurface()} indirection that resolves to whichever surface is
currently bound, while \texttt{Present()} always targets the shadow backbuffer from
\S\ref{sec:freedirect-shadow-backbuffer} regardless of what is bound at the time --- an intentional
asymmetry, since only the
shadow backbuffer is the only surface copied to the primary.

\texttt{freedirect\_texture\_rendertarget\_test.cpp} verifies this bind-redirect concretely rather than
by code inspection alone, using two distinct, deliberately different clear colors so a
readback that returned either the wrong surface's color or a stale value would fail visibly:

\begin{lstlisting}
dev.Clear(Color(20, 40, 60, 255));           // shadow backbuffer's own color, established first

auto rt = std::make_unique<RenderTarget2D>(dev, 8, 8);
dev.SetRenderTarget(rt.get());
dev.Clear(Color(210, 30, 90, 255));          // the render target's own, different color
// GetBackBufferData() now reads back (210,30,90) -- the render target's surface, not the shadow's.

dev.SetRenderTarget(nullptr);
// GetBackBufferData() over the same region now reads back (20,40,60) again --
// unbinding really redirects Clear()/ReadBackbuffer() back to Impl::backBuffer.
\end{lstlisting}

Two further checks round out the same test. \texttt{RenderTargetUsage::DiscardContents}
auto-clearing to black on rebind needed \emph{zero} FREEDIRECT-specific code at all --- it
falls straight out of shared \texttt{GraphicsDevice.cpp} logic once bind/\texttt{Clear}/read
were wired correctly, the same "comes for free" shape Chapter~\ref{ch:textures-rendertargets}
documents for other renderers. This is also a CNA/FNA identity difference: CNA repeats that
clear on an identical currently-bound target, while FNA treats the redundant call as a no-op.
The persistent DirectDraw surface makes Preserve and Platform retain color in source, but the
test proves only the Discard-to-black case; presentation usage is ignored.
\texttt{SetRenderTargets()} with two or more bindings throws, honestly:
\texttt{IDirectDrawSurface} has exactly one active surface at a time, the identical
single-active-surface conclusion Chapter~\ref{ch:sdl-renderer} reached for
\texttt{SDL\_RENDERER} for an unrelated reason (that renderer's SDL2 render target
API, not DirectDraw's own surface model) --- two structurally different renderers landing on
the same MRT answer for two different underlying reasons.

The same test confirms the size boundary: a $5000 \times 5000$ \cnaclass{Texture2D}
construction throws. \texttt{free-direct::CreateSurface} enforces a fixed
$4096 \times 4096$ cap, and the renderer propagates that exception.

\section{Public depth/stencil clear}

The ordinary colour hook locks and fills the complete active surface with all four requested
channels, and target-redirection tests prove that ``active'' changes between the shadow
backbuffer and a \cnaclass{RenderTarget2D}. Depth and stencil follow two different layers.
\cnaclass{GraphicsDevice} sees \texttt{SupportsDepthStencil()==false}, strips both flags, and
returns silently when no target bit remains --- the same depthless-target rule current FNA
uses. \texttt{freedirect\_no3d\_test.cpp} separately invokes all six renderer combination hooks and
proves their \texttt{ThrowNo3D} errors precisely because those hooks are unreachable through
the public mask. Treating either result as evidence for the other would conflate shared and
renderer contracts; \S\ref{sec:backend-clear-matrix} keeps them separate.

\section{Known limitations at the pin}

No 3D pipeline exists, matching \texttt{free-direct}'s own explicit ``Direct3D not
implemented'' stance; 8-bit palette surfaces and \texttt{GetDC} / \texttt{SetPalette}-style
GDI-adjacent calls are unsupported, since XNA itself has no palette-texture concept to begin
with; mip levels above zero throw, since \texttt{IDirectDrawSurface} has no native mip chain;
and \texttt{SetPresentationMode()} is honestly downgraded from an earlier full-pass mark to a documented
partial: it stores the requested presentation mode but never changes the actual physical
output scale and always presents in letterbox. Closing this limitation requires a change in
\texttt{free-direct}, not only CNA.

The same stored-versus-applied distinction affects two ordinary XNA properties:
\texttt{Dx3Graphics}\allowbreak\texttt{Renderer} overrides neither \texttt{SetViewport()} nor
\texttt{SetScissorRect()}, and it has no rasterizer-state hook. Both values therefore
round-trip through \cnaclass{GraphicsDevice} without changing the 2D DirectDraw/SDL output.
The SDL renderer that \texttt{free-direct} owns privately does not rescue this path: CNA never
routes these graphics-state setters into it (unlike the separate public input-discovery route
below). There is no dedicated negative pixel test, so this is a source-proven permanent 2D
boundary rather than an executed claim.

There is a related resize defect that is easier to miss because the new logical surfaces
\emph{are} created correctly. Once the first \texttt{Present()} has configured
\texttt{free-direct}'s internal renderer, a later \texttt{SetVirtualResolution()} replaces CNA's
primary and shadow surfaces but cannot reset the library's private
\texttt{logicalPresentationSet\_} flag. The next frame therefore uploads the new-sized primary
texture into the \emph{old} logical-presentation rectangle. CNA's input-transform methods
do recompute their scale from the new logical dimensions, but the public input path does not
call them: it discovers \texttt{free-direct}'s associated SDL renderer first, as the section
below explains. That renderer retains the same old logical-presentation rectangle as physical
output, so public input and presentation remain tied to the old coordinate canvas while the
game's newly allocated surfaces use the new size. Closing that three-way mismatch needs
\texttt{free-direct} to reapply logical presentation when the primary changes, not an
unilateral CNA-side approximation.

\texttt{DirectSound} / \texttt{DirectPlay} networking are explicitly out
of scope for this renderer by the project owner's own instruction --- CNA's own audio and
networking go through the ordinary renderer-agnostic paths from Chapter~\ref{ch:audio-system}
and Chapter~\ref{ch:networking} of this book, not through \texttt{free-direct} at all.

\section{The transform math is sound, but public input takes another route}

The renderer implements both \texttt{TransformWindowToLogical()} and its inverse from
the physical \texttt{SDL\_Window} size. For a logical size $(L_w,L_h)$ and physical window size
$(P_w,P_h)$, they recompute the same centered-letterbox transform that
\texttt{free-direct} uses:

\[
s = \min(P_w/L_w, P_h/L_h), \qquad
o_x = (P_w - L_w s)/2, \qquad o_y = (P_h - L_h s)/2.
\]

A logical point maps to $(x s + o_x, y s + o_y)$; a mouse/window point maps back through
$((x-o_x)/s, (y-o_y)/s)$. Thus the black bars introduced by the otherwise unavoidable
letterbox mode are accounted for explicitly, rather than being accidentally treated as game
pixels. The calculation intentionally queries the live physical window every time it is used:
the primary/shadow surfaces retain the requested logical resolution, while window managers may
choose a quite different actual window size.

Those two correct methods are nevertheless shadowed during normal public input. Both public
callers first ask \texttt{SDL\_Get\allowbreak Renderer(window)} for an associated renderer and
use SDL's own logical-coordinate conversion when one exists. Only a renderer-less window
proceeds to the renderer registry. \texttt{free-direct} creates an SDL renderer on this same
window during \texttt{SetCooperative\allowbreak Level()}, so SDL's first tier wins even though
\texttt{Dx3Graphics\allowbreak Renderer::GetRenderer\allowbreak Internal()} returns
\texttt{nullptr}. CNA can still reach the renderer through the registry; this accessor does not
return it.

This route has two time-dependent consequences. Before the first \texttt{Present()},
\texttt{free-direct} has created the renderer but has not yet called
\texttt{SDL\_SetRenderLogicalPresentation}; public input therefore passes through at window
scale. The first presentation installs the hardcoded letterbox mapping, after which SDL's own
conversion correctly follows the displayed output. A later logical resize hits the stale
\texttt{logicalPresentationSet\_} problem above: both output and public input retain the old
SDL mapping, while the new game surfaces no longer share that coordinate size.

The dedicated \texttt{Dx3\_LogicalTransform} test avoids a fragile fixed-size expectation for
that reason. It asks the running window for its physical dimensions, requests a $64\times64$
logical canvas, then verifies four properties: logical-to-window-to-logical round trips preserve
five points (corners and centre); logical $(32,32)$ becomes the physical geometric centre;
equal horizontal and vertical logical steps have equal physical length; and that measured length
is exactly $\min(P_w/64,P_h/64)$. These checks prove both the centering and the absence of
silent anisotropic stretch on the active display, including the
environment where an attempted \texttt{SDL\_SetWindowSize()} was not honoured by the window
manager. They do \emph{not} prove the public route above: the test obtains
\texttt{FreeDirectRenderer\&} and invokes both methods directly, bypassing
\texttt{Mouse::SetPosition()}, \texttt{SdlInputBridge}, and the SDL-renderer-first branch.
Chapter~\ref{ch:backend-architecture},
\S\ref{sec:backend-input-coordinate-matrix}, places that proof boundary in the full renderer
matrix.

\section{The fixed 32-bit surface decision}

\texttt{FreeDirectRenderer.cpp} creates the primary, shadow backbuffer, render targets, and
\cnaclass{Texture2D} surfaces at 32 bits per pixel. It has no palette/8-bit route. This makes two
limitations in \texttt{free-direct}'s \texttt{docs/directdraw-limitations.md} unreachable from
CNA:

\begin{description}
  \item[The discarded bit-depth parameter.] \texttt{IDirectDraw::SetDisplayMode}'s
        \texttt{dwBPP} parameter is silently discarded --- \texttt{free-direct}'s own
        \texttt{CreateSurface} always allocates a
        32bpp primary regardless of the requested depth. CNA's \texttt{CreateSurfaces()} always
        requests 32, so its request and the library result agree.
  \item[The mixed-depth blit silent-no-op.] \texttt{DirectDrawSurfaceImpl}'s \texttt{BlitFrom}
        method implements same-depth 8-on-8 and 32-on-32 copies; a mixed pair copies nothing.
        CNA creates only 32bpp surfaces, so every renderer-issued blit follows the 32-on-32 path.
\end{description}

XNA's \cnaclass{SurfaceFormat} has no 8-bit indexed-color member, so this fixed direct-color
choice fits CNA's public texture model. Direct users of \texttt{free-direct} that request 8-bit
surfaces or mixed-depth blits remain exposed to the underlying limitations.

\begin{sourcenote}
\cnaclass{PresentationParameters} cannot change that decision. The FREEDIRECT factory reads the
window, virtual size, and presentation mode, but discards
\texttt{BackBuffer\allowbreak Format}, \texttt{DepthStencil\allowbreak Format}, and
\texttt{IsFull\allowbreak Screen}. Shared device code may still ask SDL to change the window, while
free-direct remains at \texttt{DDSCL\_NORMAL}; there is no exclusive DirectDraw transition.
The primary and shadow buffers stay fixed 32-bpp, and no depth attachment exists. Thus the
public fields can round-trip while the native answer remains unchanged, exactly as the
cross-renderer matrix in \S\ref{sec:backend-presentation-format-matrix} records.
\end{sourcenote}

\subsection{Why raw factors do not identify a preset}

The blend-preset defect becomes visible with concrete factor values.
\texttt{BlendState.cpp} constructs \cnaclass{BlendState::AlphaBlend} with two color factors:
\texttt{ColorSourceBlend} is \texttt{Blend::One}, and \texttt{ColorDestinationBlend} is
\texttt{Blend::InverseSourceAlpha}. A hypothetical custom state sharing exactly those two
factors but pairing them with
\texttt{BlendFunction::Subtract} instead of the preset's own (default) \texttt{Add} computes a
different result whenever the destination pixel is not fully black:
$\mathit{src} + \mathit{dst} \times (1 - \mathit{srcA})$ versus
$\mathit{src} - \mathit{dst} \times (1 - \mathit{srcA})$. They coincide only when the
destination term is zero, as it is for \cnaclass{BlendState::Opaque}; AlphaBlend's nonzero
destination factor exposes the mismatch over existing color. A fixture that draws only onto
cleared black can therefore miss it. Preset matching now compares both
\cnaclass{BlendFunction} fields as well as the raw factors.

At the pin, the renderer plan records no open BLOCKED decisions. That plan status does not
broaden the scoped execution evidence above.
